\section{Problemática}
\label{sec:problematica}

Esta sección corresponde a la \textbf{situación real}: síntomas observables del sector ---casos,
cifras, costos--- antes de abstraer el objeto o formular el problema tecnológico (§\ref{sec:planteamiento}).

Durante años la integración bancaria dependió de acuerdos cerrados, contratos \textit{ad hoc} y
documentación dispare. El resultado fueron silos, modelos duplicados y poca reutilización de
capacidades entre líneas de negocio y socios tecnológicos.

OpenAPI y las APIs REST unificaron \textit{cómo} se describe técnicamente una interfaz (Casas et
al., 2021). Describir operaciones HTTP no garantiza, sin embargo, que la API exprese el mismo
concepto de negocio en todos los equipos.

\textbf{Casuística observable.} Casas et al.\ (2021) mapearon 47 estudios primarios: solo el
40\,\% aborda un dominio concreto; la mayoría trabaja con APIs genéricas. Esa proporción resume una
heterogeneidad de modelado que persiste pese a la adopción de OpenAPI.

Golmohammadi et al.\ (2024) revisaron 92 trabajos sobre prueba de APIs REST. Subrayan el
\textbf{costo} de validar servicios ya desplegados y la escasez de casos industriales
reproducibles.

En \textit{open banking}, la regulación británica obligó desde 2018 a nueve grandes bancos de
\textit{retail} a exponer APIs estandarizadas (CMA, 2016). Converger esas interfaces con modelos de
negocio de referencia sigue siendo, en la práctica, un ejercicio de gobernanza más que de despliegue
de especificaciones.

Como \textbf{episodio documentado}, la migración de plataformas del TSB (Reino Unido, 2018) dejó
temporalmente sin servicio a cerca de \textbf{1,9 millones de clientes}. Los costes públicos hablan
de \textbf{cientos de millones de libras} y de semanas con operación restringida (House of Commons
Treasury Committee, 2018).

El episodio no cita OpenAPI ni BIAN, pero ilustra el \textbf{dolor operativo} cuando negocio y
tecnología evolucionan sin correspondencia verificable: retrabajo, ensayo--error en adaptadores,
riesgo reputacional y regulatorio.

Farzi (2021) vincula arquitectura empresarial y modelos como BIAN con la agilidad bancaria. Advierte
barreras de madurez, heterogeneidad de implementaciones y distancia entre la arquitectura
\textit{to-be} y lo desplegado.

En el día a día, desarrollo publica contratos OpenAPI pensando en la implementación inmediata;
arquitectura y negocio conservan BIAN como referencia. La sincronización entre ambos planos sigue
siendo manual, parcial y sin métrica objetiva de coherencia.

Shabbir et al.\ (2025) señalan que la calidad industrial de las APIs depende de consistencia,
usabilidad y mantenibilidad. Esos requisitos se resienten cuando la documentación técnica no refleja
el dominio de negocio. En arquitecturas API bancarias eso se traduce en:

\begin{itemize}
  \item \textbf{Fragmentación semántica:} mismos conceptos nombrados de forma distinta entre
        OpenAPI y BIAN (p.\ ej., \texttt{PaymentOrder} frente a \texttt{paymentRequest}).
  \item \textbf{Inconsistencias estructurales:} operaciones HTTP que no corresponden a
        \textit{behaviors} del \ServiceDomain{}.
  \item \textbf{Deriva documental:} versiones de contratos OpenAPI sin trazabilidad respecto al
        modelo BIAN de referencia.
  \item \textbf{Evaluación no sistemática:} revisiones manuales o herramientas genéricas no
        adaptadas al par contrato OpenAPI bancario --- \ServiceDomain{} BIAN antes del despliegue.
\end{itemize}

Los síntomas condicionan gobernanza, interoperabilidad en ecosistemas abiertos y la confianza de
\textit{stakeholders} en que una API representa una capacidad bancaria estandarizada. Por ello el
§\ref{sec:planteamiento} pasa al diagnóstico subyacente y al problema tecnológico que ancla la
matriz (Anexo~C).
